% Source PDF: source/普通高考/2014/2014湖北理.pdf, page 1; vector flowchart.
% Grid evidence: tmp/review_2014_hubei_science/grid/page-1-grid100.jpg and
%   q13_original_final.png (no-grid crop from the source page).
% Elements: rounded start/end terminals; input parallelogram; subtraction
% rectangle b=D(a)-I(a); decision diamond b=a?; assignment rectangle a=b;
% output parallelogram; directed connectors and 是/否 labels.
% Complete node -> directed-edge list from the source:
%   start -> input -> calc -> test; test --是--> output -> finish;
%   test --否--> assign -> main entry above calc. The loop return has one
%   left-pointing arrow into the main vertical stem; the downward main arrow
%   into calc remains independent. All borders/connectors are solid.

\begin{tikzpicture}[
  >=Stealth,
  line width=0.45pt,
  line join=round,
  line cap=round,
  font=\normalsize,
  flowtext/.style={anchor=center, text centered,
    execute at begin node={\setlength{\parindent}{0pt}}},
  every node/.style={flowtext},
  flowbox/.style={flowtext, draw, minimum height=0.68cm,
    inner xsep=4pt, inner ysep=2.5pt},
  branchlabel/.style={flowtext, inner xsep=1pt, inner ysep=1pt},
  terminal/.style={flowbox, rounded corners=0.18cm, minimum width=1.45cm},
  process/.style={flowbox},
  decision/.style={flowbox, diamond, aspect=2.7, minimum width=3.45cm,
    minimum height=0.95cm},
  io/.style={flowbox, trapezium, trapezium left angle=72,
    trapezium right angle=72, minimum width=2.35cm, minimum height=0.76cm},
]
  \node[terminal] (start) at (0,9.55) {开始};
  \node[io] (input) at (0,8.05) {输入 $a$};
  \node[process, minimum width=3.75cm, minimum height=0.86cm]
    (calc) at (0,6.50) {$b=D(a)-I(a)$};
  \node[decision] (test) at (0,4.70) {$b=a?$};
  \node[process, minimum width=1.85cm] (assign) at (3.45,4.70) {$a=b$};
  \node[io] (output) at (0,3.10) {输出 $b$};
  \node[terminal] (finish) at (0,1.70) {结束};

  % Main trunk and the terminating branch.
  \draw[->] (start.south) -- (input.north);
  % Split the trunk at the loop-entry level so the return arrow merges into
  % the stem without sharing a directed segment or creating a second arrow.
  \coordinate (merge-level) at (0,7.20);
  \coordinate (merge-pre) at (0.35,7.20);
  \draw (input.south) -- (merge-level);
  \draw[->] (merge-level) -- (calc.north);
  \draw[->] (calc.south) -- (test.north);
  \draw[->] (test.south) -- node[branchlabel, left=2pt] {是} (output.north);
  \draw[->] (output.south) -- (finish.north);

  % 否 branch updates a and returns to the main stem above calc. The return
  % arrow is horizontal and points left; it does not create a second vertical
  % arrow at the calc box.
  \draw[->] (test.east) -- node[branchlabel, above=2pt] {否} (assign.west);
  \draw[->] (assign.north) -- (3.45,7.20) -- (merge-pre);
  \draw (merge-pre) -- (merge-level);
\end{tikzpicture}
